\documentclass[11pt,a4paper]{article}

\usepackage[T1]{fontenc}
\usepackage[utf8]{inputenc}
\usepackage[spanish]{babel}
\usepackage[a4paper,margin=2.8cm]{geometry}
\usepackage{microtype}
\usepackage{xurl}
\usepackage[hidelinks]{hyperref}

\hypersetup{
  pdftitle={Paraíso perdido},
  pdfauthor={Juan José Gómez Cadenas}
}

\setlength{\parindent}{1.25em}
\setlength{\parskip}{0.35em}
\raggedbottom

\title{Paraíso perdido}
\author{Juan José Gómez Cadenas}
\date{}

\begin{document}

\maketitle

%Finalmente, el meollo de la cuestión. Mi argumento es que Quenau no es diferente a Breton, a Borges o a Lorca (y creo que el texto debe dejar eso claro).Todos ellos exploran la psique --el alma--- humana, cada uno a su manera y con una variedad de técnicas, que, en todos los casos, requiere enorme trabajo, disciplina y dedicación. Un escritor que interacciona con un LLM delega "siempre" alguno de esos aspectos a la máquina y lo hace solo por incompetencia o pereza. Mi punto es un LLM no puede producir texto genuinamente humano porque carece de la psique humana y que la delegación hace que el escritor renuncia a parte de lo que le define. Mi discusión contigo demuestra que a menudo generas texto coherente y argumentos plausibles que, sin embargo son endebles o erróneos.
%
%Había leído mal el centro de tu argumento. Queneau pertenece, junto con Breton, Borges y Lorca, al lado de quienes hacen el trabajo de explorar la experiencia humana. Sus diferencias de técnica son secundarias respecto de ese compromiso. La combinatoria es uno de sus instrumentos; construirla exige el trabajo que después parece desaparecer cuando alguien mueve las tiras.
%
%El artículo lo dice al final —«todos […] tenían razón a su manera»—, pero antes la insistencia en la magia, el timo y la falta de sentido puede hacer que Queneau parezca representar precisamente lo que impugnas. Conviene distinguir con claridad el trabajo creador de Queneau de la operación de recombinar sus versos. Esa distinción permite conservar íntegra tu crítica del cuarteto sin convertirla en una descalificación de su autor.
%
%Tu tesis sobre la delegación tampoco se reduce a comparar productos. Sostienes que el esfuerzo de escribir forma parte de lo que constituye al escritor: delegarlo tiene un coste aunque el resultado sea aceptable. Y que la procedencia humana importa porque las palabras expresan una experiencia vivida, no solo porque puedan suscitar una reacción en el lector. Eso responde a mi insistencia anterior en si una máquina podría conmover: estaba cambiando el criterio de tu argumento.
%
%Sí discutiría «solo por incompetencia o pereza». Esa atribución universal de motivos no se sigue del coste que denuncias. Se puede delegar por curiosidad, presión o una decisión consciente de prioridades; todavía puedes sostener que con ello se renuncia a algo valioso.
%
%Esta conversación, además, ofrece ejemplos concretos de tu objeción: he producido distinciones plausibles apoyadas en lecturas equivocadas, y te he obligado a corregirlas. En esos momentos la fluidez verbal encubría una comprensión defectuosa. No era el trabajo intelectual que aparentaba ser.
%  Worked for 33s • 17:06
%
%Sin embargo si valoro el punto de recalcar que la creación de Quenau es tan válida como cualquier otra y también recalcar que lo que aporta el humano es, en cierto modo lo inverso de lo que ocurre en ajedrez. Los motores juegan tan bien que los equipos humano-híbrido no son realmente equipos, la máquina supera el humano por demasiado Elo. Creo que en la escritura (seria) ocurre lo mismo, pero al revés. La experiencia humana profunda hace que el humano "supere" (en esa riqueza) a la "experiencia" verbal del LLM.
%
%Esa inversión precisa tu argumento: la colaboración no garantiza una aportación equivalente de los colaboradores. En el ajedrez, la superioridad de cálculo y evaluación del motor puede volver irrelevante el criterio del humano. En la escritura seria, sostienes que la experiencia humana aporta una profundidad que la capacidad verbal del modelo no puede suplir. La facilidad para producir frases puede ocultar esa desigualdad.
%
%Queneau encaja perfectamente ahí. Su creación incluye inventar las reglas, escribir los versos y decidir qué posibilidades merece la pena abrir. Toda esa actividad pertenece a su experiencia e inteligencia como escritor. El mecanismo combinatorio trabaja sobre lo que él ha construido; no aporta una experiencia propia. La validez de su obra no depende de que cada combinación resulte un poema memorable.
%
%La diferencia que conviene conservar es que en ajedrez podemos medir la superioridad por los resultados; en literatura tienes que mostrar qué riqueza valoras. Eso es precisamente lo que hacen Lorca, Virgilio y Rilke en tu artículo: permiten al lector reconocer y juzgar esa riqueza. No necesitas inventarle un Elo literario.
%
%Con esta formulación, el problema de delegar se vuelve más concreto: el escritor puede dejar que la soltura verbal del modelo decida por él antes de haber encontrado una expresión nacida de su propia experiencia. El texto sale antes, pero parte del trabajo mediante el cual habría descubierto qué tenía que decir queda sin hacer.
%  Worked for 24s • 17:10
%
%Dos comentarios aquí. Una, el ELO en literatura es intersubjetivo, pero no por eso menos potente. Magnus no tiene más ELO que Rilke, aunque los midamos de forma muy distinta. Por otra parte, dejar que la soltura vebal del modelo ayude al escritor es exactamente lo contrario que proponen tanto Breton (que le pide al escritor que explore su psique mediante escritura automática) como Queneau (que le pide que se ciña a unas reglas muy estrictas). Delegar al modelo es pasar de ambos.
%
%Mi distinción entre las dos medidas no debía sugerir que el juicio literario tiene menos fuerza. La grandeza de Rilke no queda debilitada porque se reconozca mediante una valoración intersubjetiva y no mediante una puntuación competitiva. La analogía que propones compara superioridades dentro de cada actividad.
%
%Y tu segundo comentario precisa la contradicción de apelar al Oulipo para justificar la delegación. Breton exige que el escritor se exponga a sus propias asociaciones; Queneau exige que encuentre soluciones bajo restricciones elegidas. Aunque los procedimientos difieran, el trabajo que imponen al escritor es constitutivo de la creación. Encargarle al modelo ese trabajo no equivale a practicar ninguna de esas técnicas.
%
%Por eso mi expresión «que la soltura verbal del modelo ayude» concedía de antemano lo que estás discutiendo: que facilitar la producción de frases ayude necesariamente a escribir. En tu argumento, puede eliminar precisamente la dificultad que obliga al escritor a explorar, descubrir y transformar su experiencia. Lo que Ángel presenta como asistencia tú lo cuestionas como sustitución del trabajo creador.

Hace ahora poco más de un siglo, David Hilbert acuñaba la frase: «Nadie podrá expulsarnos del paraíso que Cantor ha creado para nosotros». Estamos en Hamburgo, durante la conferencia «\emph{Über das Unendliche}»; la fecha es 1925 y Hilbert lleva un cuarto de siglo embarcado en una de las empresas intelectuales más ambiciosas del \emph{novecento}. Quiere rescatar a las matemáticas de las paradojas de la teoría de conjuntos, formuladas por Bertrand Russell y otros entre finales del siglo XIX y principios del XX. Buena parte de esas paradojas tienen su origen en el trabajo de Cantor, quien había descubierto la existencia de un extraño Paraíso, poblado por jerarquías de infinitos ---los llamados cardinales transfinitos--- que los bienaventurados podían manipular. Hilbert tiene como objetivo eliminar las paradojas que amenazan con expulsar a los matemáticos de esa Tierra Prometida. Para ello, propone formalizar por completo las matemáticas y demostrar que un sistema así construido no puede producir contradicciones. Su programa implica desarrollar una metodología que permita demostrar teoremas sin hacer otra cosa, en última instancia, que manipular símbolos. 

Es decir, Hilbert quiere construir una máquina de fabricar teoremas.

Unos pocos años más tarde, Alan Turing reformula la idea en términos de un ingenio capaz de leer instrucciones codificadas en una cinta sin fin ---papel perforado en su época--- y escribir sus resultados en otra cinta análoga. Ese artefacto puede codificar cualquier algoritmo y, por tanto, ejecutar las secuencias de manipulaciones simbólicas que llevan a la demostración de cualquier teorema. De esta manera, la entelequia de Hilbert se transforma en un objeto todavía ideal pero mucho más cerca de lo real. La máquina de Turing. 

Permíteme, incisiva lectora, atento lector, un pequeño \emph{cliffhanger} antes de continuar con esta historia. Saltemos a la noche del 24 de noviembre de 1960, a una cena que podemos imaginar desenfadada y solemne a partes iguales. Una cena bien regada con borgoñas y burdeos, donde no faltan ni el \emph{confit} de pato ni una selección de quesos bretones, rematada con buen \emph{brandy} y cigarros cubanos. El anfitrión es el célebre literato y matemático aficionado Raymond Queneau. Hay una decena de comensales, que incluyen la flor y nata de la \emph{intelligentsia divine} francesa. Entre los invitados encontramos científicos, filósofos, ajedrecistas, no pocos matemáticos ---incluyendo a Claude Berge, uno de los padres de la teoría de grafos--- y escritores notables, como Noël Arnaud y Jacques Duchateau. Todos tienen ganas de divertirse y de liarla un poco, sobre todo una vez que se han vaciado las tres primeras botellas de calvados. En algún momento de la larga velada, acuerdan formar el Seminario de literatura experimental, que acaba por convertirse en el Oulipo (\emph{Ouvroir de littérature potentielle}, o «taller de literatura potencial»). Los miembros fundadores se comprometen a serlo \emph{ad infinitum}, incluyendo continuidad \emph{post mortem} ---se acuerda que los que pase a mejor vida podrán eludir participar en futuras reuniones excusándose por defunción---. La broma podría haberse quedado ahí, pero el taller en cuestión se hace famoso y acaba por enrolar a luminarias como Georges Perec, Italo Calvino y Marcel Duchamp.

Los miembros del Oulipo postulan que la literatura, como las matemáticas, puede, en último término, reducirse a un algoritmo capaz de manejar símbolos ---los del lenguaje--- a partir de una serie de reglas fijas. La idea recuerda mucho los fundamentos del programa de Hilbert y el propio Queneau la desarrolla en uno de sus textos, titulado nada menos que «\emph{Les fondements de la littérature d'après David Hilbert}». Queneau toma los axiomas de los \emph{Grundlagen der Geometrie} y sustituye puntos, rectas y planos por palabras, frases y párrafos. Es decir, literalmente, construye teoremas lingüístico-literarios.

Permíteme, curiosa lectora, siemprealerta lector, un poco de contexto. Queneau y los oulipianos se oponen frontalmente al movimiento surrealista de Breton del que el propio Queneau ha formado parte en su juventud. El Oulipo puede explicarse, al menos parcialmente, como la reacción inmunitaria de parte de la élite literaria francesa al reinado ---que ya se iba convirtiendo en tiranía--- de Breton y su corte. Así, donde los unos predican la necesidad de la escritura automática como método para explorar el subconsciente y quizás aprehender lo inefable, los otros se burlan de una técnica muy proclive a producir lo que hoy no dudaríamos en calificar de \emph{slop} y proponen a cambio encadenar al literato a un conjunto de reglas tan estricto que la creatividad no tenga otro remedio que surgir para producir algo respetándolas\footnote{Una propuesta que evoca el célebre chiste en el que un explorador se pasea por la jungla acarreando un pesado yunque, de tal forma que, en caso de encontrarse con un león, puede tirarlo y corre más rápido.}.

%Ciertamente, los que pasamos por el sarampión surrealista podemos simpatizar a ratos con los oulipianos. Pero permíteme, atenta lectora, sagaz lector, que te ofrezca estas líneas (que cito de memoria, leídas hace cuarenta años y nunca olvidadas):
%
%\begin{verse}
%Las criaturas de la luna huelen y rondan las cabañas,\\*
%vendrán las iguanas vivas a morder a los hombres que no sueñan\\*
%y el que huye con el corazón destrozado, encontrará por las esquinas\\*
%al increíble cocodrilo, quieto, bajo la tierna protesta de los astros.
%\end{verse}
%
%Se ha hablado hasta la saciedad de los gitanos de Lorca, pero esas criaturas de la luna, esas iguanas vivas, ese impasible cocodrilo quieto bajo los astros, todavía me ponen la carne de gallina, todavía me hacen estremecer, todavía me parece que capturan algo inefable, inasible, profundamente misterioso del alma humana.

Hacia 1960, Queneau publica sus célebres
\emph{Cent mille milliards de poèmes}. Se trata de diez sonetos con las mismas rimas, en los que cada verso se escribe en una tira de papel ofreciendo al lector $10^{14}$ posibles combinaciones. Es, como bien señala mi querido Ángel Fernández en su reciente artículo «Las tijeras de Queneau en la Biblioteca de Babel», una deliciosa vuelta de tuerca a la noción de la biblioteca de Babel. Donde Borges solo puede ofrecernos horror combinatorio, Queneau nos ofrece un número gigantesco de posibles poemas, todos los cuales tendrían, al menos formalmente, sentido.

Queneau acaba de construir algo similar a la máquina de Turing. Una máquina de hacer sonetos, capaz de producir un número astronómico de ellos siguiendo un simple algoritmo combinatorio. Por ejemplo: 

%\emph{Al menos formalmente}. El propio Ángel nos da un ejemplo:
%
%
%\begin{quotation}
%«\emph{Le roi de la pampa retourne sa chemise / depuis que lord Elgin négligea ses naseaux}»
%
%«El rey de la pampa cambia de camisa / desde que lord Elgin descuidó sus ollares» no parece el comienzo de un poema de Neruda. La combinación de personajes parece arbitraria en la forma en que trastoca y construye las palabras y, aunque es clara como una escritura notarial, no parece llevarnos a la exaltación del espíritu que uno espera de la poesía.
%\end{quotation}
%
%Anota, precisa lectora, sofisticado lector, este detalle: \emph{no parece llevarnos a la exaltación del espíritu que uno espera de la poesía.} Anota también el hecho de que la coherencia formal no implica sentido. En la biblioteca de Babel, la mayoría de los libros serían combinaciones arbitrarias como esta:
%
%\begin{verse}
%\ttfamily
%hmarusfaqictruyqsurabnddcpsvgsmomzpundh\\*
%jjzdqvnphoujxtapxñiundgnhqetdimnqcicemuazdmughviñ\\*
%yoiveqsceifphdefrhhndvypgmyxñxtzñuijox\\*
%ñbfbcqbtmgxeajpinqntmgccuuvjmuatjvxgbqezossz
%\end{verse}
%
%Pero una variación cualquiera de los versos de Queneau se leería así:

\begin{verse}
\itshape
Il se penche il voudrait attraper sa valise\\*
pour la mettre à sécher aux cornes des taureaux\\*
une toge il portait qui n’était pas de mise\\*
il chantait tout de même oui mais il chanta faux
\end{verse}

Traducido al español:

\begin{verse}
Se inclina, quisiera alcanzar su maleta\\*
para ponerla a secar en los cuernos de los toros;\\*
llevaba una toga que no venía al caso,\\*
cantaba de todos modos, sí, pero desafinaba.
\end{verse}

¿Tienen sentido los versos o solo lo parece? Los seres humanos le buscamos sentido a todo y más todavía cuando se trata del lenguaje, así que nos esforzamos por imaginar a un señor vestido con una toga algo desaliñada, cantando desafinadamente mientras se esfuerza en alcanzar su maleta (mojada) para ponerla a secar nada menos que en los cuernos de los toros. Por una parte, el cuarteto, no dice nada. Por otra parte, es una pequeña descarga eléctrica. Y nada obliga al lector a escoger los versos al azar. Podemos imaginar que este se torna escritor mediante el acto de seleccionar cuidadosamente los versos que le permitan expresarse, o inquietarnos, o explorar extraños nuevos territorios, que pueden incluir maletas que cuelgan en un enorme tendedero hecho de asta de miuras. 

Sostiene Ángel Fernández que la escritura híbrida entre hombre y máquina es posible, utilizando un sistema similar al de los cien billones de poemas y la propuesta, a primera vista, es muy sugestiva. Uno puede imaginar al LLM como la máquina de producir versos en tiras de papel ---una máquina de Queneau--- y al escritor como aquel que los selecciona. 

Pero creo que, atractiva y todo, la analogía no se sostiene. Y para apoyar mis argumentos, quiero referirme 

No me queda claro si defiende que sea la máquina quien genere el texto y el humano el que lo ordene, si se trata de lo contrario, o si utiliza toda la metáfora oulipiana como una justificación del argumento que se oye cada vez más a menudo: «yo le digo mis ideas a la máquina y la máquina me las escribe, luego yo las repaso».

La metáfora, en ese caso, sería muy engañosa. Queneau y compañía hacen exactamente lo contrario de lo que los «coescritores» postulan. Los oulipianos trabajan denodadamente para construir los versos que respetan las arduas reglas que se han impuesto, sin delegar ese arduo trabajo ---que es el puntal en el que se apoya la creatividad--- a una máquina, igual que los surrealistas de Breton no le pedían a ChatGPT que escribiera texto al azar como sustituto de su escritura automática.

«Yo digo mis ideas y la máquina me las escribe» es una falacia cósmica. Las ideas no se formulan de un golpe, ni mucho menos. La primera vez que nos cruzan por la cabeza son poco más que una intuición, una pista emborronada de algo que no toma forma hasta que lo escribimos. El primer borrador suele ser un batiburrillo de confusas ideas que solo sirve para crear un segundo borrador. Quizás, a partir del tercero ---este artículo ya va por su quinta versión--- empezamos a entender lo que decimos. Delegar un \emph{prompt} a una máquina, digamos: «escribe un artículo que sea un espejo de “Borges en su biblioteca”, utilizando como contraargumento el Oulipo» y luego repasar lo que la máquina produce puede resultar en un texto simpático, provocador y en último término vacío, porque carece de la reflexión que el humano ha delegado y la máquina no puede aportar. Un LLM funciona por asociación ---una asociación compleja y nada trivial, desde luego--- que permite, gracias a su entrenamiento dirigido, conducir su mecanismo básico ---adivinar la palabra más probable que sigue dentro de su ventana de contexto--- a generar texto razonablemente convincente, pero no especialmente profundo ---a menudo, espectacularmente trivial---.

\emph{For the record}, nada tengo contra la coautoría hombre-máquina. Solo reclamo que se le dé a César lo que es de César y a Claude lo que es de Claude. El problema de los escritores pillados con las manos en la masa ---o debería decir en la IA--- en recientes concursos literarios es que han hecho trampas, saltándose las reglas de participación en esos concursos. Me parece excelente, como proponen algunos escritores, que se convoquen concursos híbridos y se publiquen novelas en las que Claude sea coautor(a). Les auguro, eso sí, el mismo éxito que han tenido los torneos híbridos de ajedrez. En ese bello juego-deporte que tanto me apasiona ---y en el que tan torpe soy---, se siguen con devoción las partidas entre humanos y se hace poco o ningún caso a las competiciones entre máquinas o entre equipos hombre-máquina. Extrañamente, a pesar de que las máquinas juegan mejor que nosotros, no nos interesan. Y la razón que lo explica, creo, es que el ajedrez no es solo un juego. Es un combate de boxeo psicológico, un drama silencioso, un teatro de kabuki, una batalla mortal en la que los contendientes no se levantan de la silla, una exhibición de belleza, un catálogo de errores, una explosión de intuición, un torrente de emociones reprimidas. AlphaZero puede jugar todo lo bien que quiera, pero no nos emociona. Dudo mucho que las novelas o los poemas híbridos (o puramente IA) que tanto defienden mis amigos funcionen mejor. AlphaZero juega mejor que un humano porque el ajedrez se puede reducir a un algoritmo programable en una máquina. Pero la poesía no es un algoritmo y quien crea que un LLM puede capturar las \emph{Elegías de Duino}, o no ha leído a Rilke, o no ha leído las mediocridades que producen los \emph{chatbots}.

También, para que conste, la policía IA no se inventó en nuestro zoológico de plumíferos, sino en los clubs \emph{online} de ajedrez. Y se inventó como remedio desesperado para que el ajedrez \emph{online} sobreviviera. Y solo ha tenido éxito a medias.

Me explico. No hay cosa más absurda que jugar una partida \emph{online} ayudándose de un motor. La gracia del ajedrez es, precisamente, jugar. Cuando un tramposo deja que el robot mueva por él las piezas, le está haciendo perder el tiempo a su oponente y se está engañando a sí mismo en la misma medida que se engañaría un halterófilo que hiciera sentadillas con 500 kilos..., ayudado por una máquina que se los levantara. Y aun así, el ajedrez en la red está infectado por los tramposos. Es una epidemia imparable e incurable. En los clubs \emph{online} la policía IA ---consistente en programas que detectan a los que hacen trampas fijándose en que juegan con la perfección de una máquina--- consigue mitigar algo esa pandemia, pero no lo suficiente. Hay jugadores ---sobre todo entre los que echan partidas largas o por turnos--- que tienen la paciencia ---y la deshonestidad--- de equivocarse a propósito de vez en cuando, elegir sistemáticamente la segunda o tercera mejor jugada que sugiere la máquina en lugar de la primera, todo para estafar al detector y ganar la partida... ¿con qué propósito?

Y sin embargo, los clubs \emph{online} siguen utilizando detectores IA con el único propósito de \emph{desanimar a la mayoría de los tramposos}, esto es, para mitigar el \emph{slop}. Lo consiguen solo a medias. Hoy en día, todos los que jugamos \emph{online} nos limitamos a partidas cortas. Si queremos echar una partida larga, nos vamos al club. El ajedrez por correspondencia, a todos los efectos, es cosa del pasado.

Mi argumento para favorecer Pangram no es que piense que se puede detectar a todos los tramposos ---y repito, llamo tramposo al que escribe con IA sin declararlo---, sino desanimar a los ejércitos de narcisos armados de un \emph{chatbot}. Como en el caso del ajedrez \emph{online}, creo que va a ser muy difícil, pero insisto en ello porque creo que el día que se abran las compuertas a todo el mundo que sabe pedir el \emph{Quijote} con un \emph{prompt}, los pocos que todavía leen dejarán de hacerlo.

Pero volvamos a Hilbert. Unos años después de su célebre declaración exigiendo el Paraíso de Cantor, Gödel publica sus teoremas de incompletitud que pueden resumirse así: «todo sistema formal contiene teoremas ciertos que no pueden demostrarse». Poco después, el mismo Turing da una versión informática con su célebre \emph{halting problem}, demostrando que no existe un algoritmo universal.

¿Debemos entender que Gödel y Turing expulsaron a los matemáticos del Paraíso? Yo creo que no. Creo que los teoremas de incompletitud revelan que la realidad ---y por tanto la mente humana, ya que la realidad pasa por ella--- contiene una parte irreduciblemente no algorítmica. Y en esa irreducibilidad fuera del alcance de un algoritmo ---y de una máquina--- está nuestra última Tierra Prometida.

Mientras mi amigo Ángel y yo nos entretenemos con juegos florales, el mundo de las matemáticas está patas arriba. La avalancha de teoremas demostrados con IA ha supuesto un curioso cambio de guion. De repente, el sueño de Hilbert se hace parcialmente cierto ---es posible demostrar teorema tras teorema de forma mecánica, simplemente mediante una máquina que manipula símbolos--- y todos los matemáticos del mundo se rasgan la camisa. Terence Tao, una de las más icónicas luminarias del momento, insiste en que la demostración automática, lejos de llevarnos al Paraíso, nos conduce al Averno. Alega que sin recorrer el largo camino a Ítaca, no hay comprensión y sin comprensión no hay nada. Es lo mismo que yo defiendo. Un texto literario surge de una extraña combinación de inspiración, lucidez y trabajo. Todos los pedantes que penan en el infierno de la Biblioteca (Borges, Queneau, Breton) tenían razón a su manera, porque cada uno de ellos propuso una manera de acercarse a la elusiva verdad. Y ninguna de esas maneras consistía en lanzarle un \emph{prompt} a una máquina para que pensara por ellos.

He prometido al principio del artículo volver sobre una línea del artículo de Ángel que dice:


\begin{quotation}
y, aunque es clara como una escritura notarial, no parece llevarnos a la exaltación del espíritu que uno espera de la poesía.
\end{quotation}

Exactamente. Las combinaciones oulipianas, no digamos ya las que produce un LLM, merecen justamente el epíteto notarial, pero no se cuentan entre las cosas que merecen lágrimas.

Y las cosas que merecen lágrimas son la justificación última de la literatura, al menos para los huérfanos de Dios.
Merecen lágrimas la muerte de Dido y la de Julieta, la escena en la que la Maga descubre que su bebé ha dejado de respirar, el páramo que atraviesa un personaje de \emph{Cien años de soledad} donde el eco repite los pensamientos, Príamo besando las manos de Aquiles, el momento en el que Orfeo mira hacia atrás y pierde para siempre a Eurídice, los versos de Dylan Thomas, «\emph{Rage, rage against the dying of the light}», la tenacidad del escribiente que decide no hacerlo, la locura de Ahab, el amor desesperado del padre que lleva una sola bala en su revólver para matar a su hijo antes que abandonarlo a las alimañas que pueblan la Tierra.

Sostengo que ninguna IA ---no las que conocemos--- escribirá líneas comparables a esas. Una cosa es divertirse reordenando renglones escritos en tiras de papel y otra cosa es captar la extraña cualidad del alma humana, esa esencia irreducible en la que Gödel, Borges y Queneau se dan la mano.

¿De dónde sale esa terca inefabilidad? Precisamente del hecho de que somos seres de carne y hueso, de piel que ha sentido otra piel rozándonos, de ojos que se han maravillado contemplando el océano, de brazos que han sostenido a nuestros bebés y manos que han acariciado a nuestros amantes; somos sacos de protoplasma con un kilo y medio de materia gris en el cráneo y los intestinos llenos de mierda, y es la suma de esa mierda y esos sesos donde se conectan cien mil millones de neuronas y esa piel que tanto recuerda y esa boca que ha besado y mordido y ese deseo que nos enloquece y nos eleva la que nos convierte en seres espirituales, ángeles cargados de cadenas, capaces de entender, \emph{de intuir} la belleza de lo efímero y lo transcendente, como ningún LLM ---armado solo de palabras--- podrá jamás entender:

\begin{verse}
Pues la belleza no es otra cosa que el principio del terror,\\
y nos sorprende cuando, serenamente,\\
desdeña aniquilarnos.\\
Todo ángel es terrible.
\end{verse}

\end{document}
